\item \textbf{Cálculo del campo magnético interno en la radiación de 21 cm.}
Trabaja en la NASA analizando la radiación interestelar de $21\text{ cm}$ emitida por el hidrógeno atómico. Esta radiación proviene de una división hiperfina del estado fundamental debido al acoplamiento de espín del electrón con el espín del protón (que genera un campo magnético interno). Sabiendo que la longitud de onda de emisión es exactamente de $21.1\text{ cm}$, estime la magnitud promedio del campo magnético interno $B$ en el que reside el electrón.